\chapterbackmatter{Supplementary Exercises}

\begin{enumerate}
  \SupplementaryExercise{1}{Discrete Symmetries and Dangerous MSSM Operators}
    {(Inspired by Chung et al., ``The Soft Supersymmetry-Breaking Lagrangian: Theory and Applications'', Sections 2.3 and 10.2; independently formulated.)}
    {ch28-s01-discrete-symmetries-and-dangerous-mssm-operators}
  \begin{enumerate}
    \item \emph{Matter-Parity Test of the Renormalizable Superpotential.}
    Assign matter parity \(-1\) to every quark and lepton chiral superfield
      and \(+1\) to both Higgs superfields.  Classify all gauge-invariant
      renormalizable superpotential monomials and identify which violate
      baryon or lepton number.
    \item \emph{Proton Decay from Simultaneous RPV Couplings.}
    Combine one \(LQ\bar D\) coupling with one
      \(\bar U\bar D\bar D\) coupling and integrate out the appropriate
      right-handed down squark.  Write a representative four-fermion operator
      and derive its coefficient including flavor labels.
    \item \emph{Baryon Triality Check for Dimension-Four Operators.}
    Construct a multiplicative \(\mathbb Z_3\) charge assignment that
      permits the ordinary Yukawa and \(\mu\) terms and the lepton-number
      violating renormalizable monomials, but forbids
      \(\bar U\bar D\bar D\).  Check every monomial explicitly.
  \end{enumerate}

  \SupplementaryExercise{2}{One-Loop Unification and Superpartner Thresholds}
    {(Adapted from University of Cambridge, 2024 Part III Paper 307, \emph{Supersymmetry}, Question 3.)}
    {ch28-s02-one-loop-unification-and-superpartner-thresholds}
  \begin{enumerate}
    \item \emph{One-Loop Supersymmetric Unification Coefficients.}
    Starting from the gauge-boson, Weyl-fermion, and complex-scalar
      contributions to a one-loop beta function, derive the three MSSM
      coefficients for three generations and two Higgs doublets.  Verify the
      relative hypercharge normalization used at unification.
    \item \emph{Threshold Split Between Scalar and Fermion Partners.}
    A chiral multiplet has a light complex scalar of mass \(m_s\) and a
      heavier Weyl fermion of mass \(m_f\).  Derive its piecewise contribution
      to \(1/g_i^2(\mu)\) and express the net threshold correction in terms of
      \(\ln(m_f/m_s)\).
  \end{enumerate}

  \SupplementaryExercise{3}{Gaugino Mass Ratios at a Common Scale}
    {(Inspired by Stephen P. Martin, \emph{A Supersymmetry Primer}, version 7, Sections 6.5 and 8.1; independently formulated.)}
    {ch28-s03-gaugino-mass-ratios-at-a-common}
In minimal gauge mediation with complete unified messenger
  multiplets, prove that \(M_i/g_i^2\) is common at the messenger scale
  and remains one-loop RG invariant.  Estimate the bino, wino, and gluino
  mass ordering at an electroweak scale.

  \SupplementaryExercise{4}{Soft Flavor Tensors and Mass-Insertion Bounds}
    {(Inspired by Chung et al., ``The Soft Supersymmetry-Breaking Lagrangian: Theory and Applications'', Sections 5.1--5.2; independently formulated.)}
    {ch28-s04-soft-flavor-tensors-and-mass-insertion}
  \begin{enumerate}
    \item \emph{Soft-Term Flavor Counting in a Super-CKM Basis.}
    Diagonalize the quark Yukawa matrices and enumerate the independent
      flavor structures that remain in the left- and right-squark squared-mass
      matrices and trilinear couplings.  Explain why gauge symmetry alone does
      not align them with the quark masses.
    \item \emph{A Neutral-Kaon Mass-Insertion Bound.}
    Define \((\delta^d_{12})_{LL}=(\Delta M_{\widetilde d}^2)_{12}/
      \widetilde M^2\).  Use a gluino box and dimensional analysis to show
      that its \(\Delta S=2\) coefficient scales quadratically in this
      insertion, and compare it with the standard-model charm contribution.
    \item \emph{A Slepton-Mixing Bound from Radiative Muon Decay.}
    Insert one off-diagonal slepton squared mass into a neutralino--slepton
      loop for \(\mu\to e\gamma\).  Determine the dependence of the dipole
      coefficient and branching ratio on the insertion, electroweak coupling,
      and common superpartner mass.
  \end{enumerate}

  \SupplementaryExercise{5}{Physical Soft Phases and Dipole Decoupling}
    {(Inspired by Chung et al., ``The Soft Supersymmetry-Breaking Lagrangian: Theory and Applications'', Sections 5.1--5.2; independently formulated.)}
    {ch28-s05-physical-soft-phases-and-dipole-decoupling}
  \begin{enumerate}
    \item \emph{Rephasing-Invariant Soft CP Phases.}
    In a one-generation MSSM with complex \(\mu\), \(B\), a common gaugino
      mass \(M_{1/2}\), and a common trilinear parameter \(A\), use field
      rephasings to construct two independent CP-odd invariant phase
      combinations.
    \item \emph{Parametric Decoupling of a Supersymmetric EDM.}
    Estimate a light-fermion electric dipole moment generated by one gaugino,
      two sfermion propagators, and one left--right mass insertion.  Show
      explicitly how the answer decouples when every soft mass is scaled by a
      common factor.
  \end{enumerate}

  \SupplementaryExercise{6}{A General Two-Higgs-Doublet Sector}
    {(Adapted from David Tong, \emph{The Standard Model Example Sheet 4}, Question 6.)}
    {ch28-s06-two-higgs-superfields-and-electroweak-breaking}
  \begin{enumerate}
    \item \emph{Charges.}
    Besides the standard Higgs \(H\sim(\mathbf2,\mathbf1)_{+1/2}\), introduce
      \(\phi\sim(\mathbf2,\mathbf1)_{-1/2}\).  Determine the electric charge
      of each component of \(\phi\).
    \item \emph{General Renormalizable Potential.}
    Write every gauge-invariant Hermitian operator of dimension at most four
      built from \(H\) and \(\phi\).  Count the independent real parameters,
      allowing all coefficients permitted by gauge symmetry to be complex.
    \item \emph{Vacuum and Gauge Spectrum.}
    Suppose
      \[
        \langle H\rangle=\frac1{\sqrt2}\binom0v,\qquad
        \langle\phi\rangle=\frac1{\sqrt2}\binom{u+iw}0,
        \qquad u,v,w\in\mathbb R.
      \]
      Decide whether electromagnetism is broken.  Compute the masses of
      \(W^\pm\), \(Z\), and the photon, and identify the electric charges of
      the remaining physical scalar fields.
    \item \emph{Yukawa Couplings.}
    List all gauge-invariant quark and charged-lepton Yukawa couplings to the
      two doublets.  Explain why gauge symmetry alone permits both doublets to
      couple to each fermion species.
  \end{enumerate}

  \SupplementaryExercise{7}{Goldstone, Pseudoscalar, and Charged-Higgs Directions}
    {(Inspired by Manuel Drees, \emph{An Introduction to Supersymmetry}, Sections 4b--4c; independently formulated.)}
    {ch28-s07-goldstone-pseudoscalar-and-charged-higgs-directions}
  \begin{enumerate}
    \item \emph{Neutral Goldstone and Pseudoscalar Directions.}
    Diagonalize the imaginary neutral-scalar mass matrix in
      Eq.~\eqref{eq:28.5.26}.  Give normalized combinations for the neutral
      Goldstone boson and physical pseudoscalar, and verify their squared
      masses without solving a general quadratic equation.
    \item \emph{The Charged-Higgs Tree-Level Sum Rule.}
    Starting from the charged scalar matrix
      Eq.~\eqref{eq:28.5.34}, find the charged Goldstone direction and derive
      \(m_{H^\pm}^2=m_A^2+m_W^2\).  Explain which part of the result follows
      specifically from supersymmetric \(D\)-term quartics.
  \end{enumerate}

  \SupplementaryExercise{8}{The CP-Even Higgs Sector and Its Leading Correction}
    {(Inspired by Manuel Drees, \emph{An Introduction to Supersymmetry}, Sections 4b--4c; independently formulated.)}
    {ch28-s08-the-cp-even-higgs-sector-and}
  \begin{enumerate}
    \item \emph{CP-Even Higgs Trace, Determinant, and Bounds.}
    Compute the trace and determinant of the tree-level CP-even matrix
      Eq.~\eqref{eq:28.5.27}.  Use them to recover the two eigenvalues and
      prove \(m_h\leq\min(m_A,m_Z)\) without assuming a special value of
      \(\beta\).
    \item \emph{Light-Higgs Alignment in the Decoupling Limit.}
    Expand the CP-even mixing angle and eigenvalues through first nonzero
      order in \(m_Z^2/m_A^2\).  Show that the light scalar aligns with the
      vacuum direction and quantify the leading departure of its Yukawa
      couplings from standard-model values.
    \item \emph{Leading Stop Logarithm in the Higgs Mass.}
    Using the one-loop effective potential for a top--stop multiplet with
      negligible stop mixing, derive the leading logarithmic correction to
      the light CP-even Higgs squared mass in the large-\(m_A\) limit.  State
      why the correction is positive.
  \end{enumerate}

  \SupplementaryExercise{9}{Chargino Singular Values and Limiting Spectra}
    {(Inspired by Manuel Drees, \emph{An Introduction to Supersymmetry}, Section 4c; independently formulated.)}
    {ch28-s09-chargino-singular-values-and-limiting-spectra}
Use the trace and determinant of
  \(\mathcal M_{\mathrm C}^\dagger\mathcal M_{\mathrm C}\) for
  Eq.~\eqref{eq:28.5.51} to derive both chargino squared masses.  Then
  describe the wino-like and higgsino-like limits and the exceptional
  near-degenerate case.

  \SupplementaryExercise{10}{Neutralino and Sfermion Mixing}
    {(Inspired by Stephen P. Martin, \emph{A Supersymmetry Primer}, version 7, Sections 7.7 and 8.2--8.4; independently formulated.)}
    {ch28-s10-neutralino-and-sfermion-mixing}
  \begin{enumerate}
    \item \emph{Neutralino Mixing in the Heavy-Higgsino Limit.}
    Write the neutralino mass matrix in the bino, neutral wino, and two
      neutral higgsino basis.  For \(|\mu|\gg|M_1|,|M_2|,m_Z\), integrate out
      the higgsinos and find the leading corrections to the two gaugino-like
      masses.
    \item \emph{Sfermion Left--Right Mixing and Fermion Mass.}
    Derive the off-diagonal entry of a one-generation sfermion squared-mass
      matrix from the \(\mu\) term and the aligned trilinear soft term.  Show
      why the mixing is proportional to the corresponding fermion mass and
      compare stops with first-generation squarks.
  \end{enumerate}

  \SupplementaryExercise{11}{Messenger Stability and One-Loop Gaugino Masses}
    {(Inspired by Giudice and Rattazzi, ``Theories with Gauge-Mediated Supersymmetry Breaking'', Sections 2.1--2.3; independently formulated.)}
    {ch28-s11-messenger-stability-and-one-loop-gaugino}
  \begin{enumerate}
    \item \emph{Messenger Supertrace and the No-Tachyon Condition.}
    For \(W=\lambda S\Phi\bar\Phi\) with
      \(S=\mathcal S+\theta^2\mathcal F\), diagonalize the messenger boson and
      fermion mass matrices.  Verify the messenger supertrace and derive the
      exact inequality that prevents a tachyonic scalar.
    \item \emph{Nonidentical Spurions and One-Loop Gaugino Masses.}
    Let each messenger pair couple to a different spurion \(S_n\).  Derive
      the leading gaugino masses for all three standard-model gauge factors,
      keeping the Dynkin indices and complex ratios
      \(\mathcal F_n/\mathcal S_n\) explicit.
  \end{enumerate}

  \SupplementaryExercise{12}{Scalar Soft Masses and Unified Messenger Indices}
    {(Inspired by Giudice and Rattazzi, ``Theories with Gauge-Mediated Supersymmetry Breaking'', Sections 2.3--2.4; independently formulated.)}
    {ch28-s12-scalar-soft-masses-and-unified-messenger}
  \begin{enumerate}
    \item \emph{Two-Loop Scalar Masses from Quadratic Casimirs.}
    Starting from gauge invariance and two-loop power counting, write the
      minimal gauge-mediated squared mass of a visible scalar in terms of its
      three quadratic Casimirs.  Apply the result to \(Q\), \(\bar U\),
      \(L\), and \(\bar E\).
    \item \emph{Complete Messenger Multiplets and Unification.}
    Decompose one \(\boldsymbol5+\bar{\boldsymbol5}\) of \(SU(5)\) under the
      standard-model gauge group.  Compute its three one-loop messenger
      indices and prove that it shifts the inverse unified couplings equally
      in grand-unified normalization.
  \end{enumerate}

  \SupplementaryExercise{13}{Flavor Universality and the Gravitino LSP}
    {(Inspired by Giudice and Rattazzi, ``Theories with Gauge-Mediated Supersymmetry Breaking'', Sections 3.1--3.2; independently formulated.)}
    {ch28-s13-flavor-universality-and-the-gravitino-lsp}
  \begin{enumerate}
    \item \emph{Flavor Universality at the Messenger Scale.}
    Show that gauge-mediated squark and slepton masses are proportional to
      the identity in generation space at the messenger scale.  Identify the
      Yukawa-dependent renormalization effects that spoil exact degeneracy at
      lower scales and explain their flavor alignment.
    \item \emph{Gravitino Scaling and the Identity of the LSP.}
    Compare \(m_{3/2}\sim F/(\sqrt3M_{\rm Pl})\) with a gauge-mediated soft
      mass \(m_{\rm soft}\sim(\alpha/4\pi)F/S\).  Determine the parametric
      region in which the gravitino is the LSP and derive the
      \(F\)-dependence of an NLSP decay length.
  \end{enumerate}

  \SupplementaryExercise{14}{The Gauge-Mediation \(\mu\)--\(B\mu\) Problem}
    {(Inspired by Giudice and Rattazzi, ``Theories with Gauge-Mediated Supersymmetry Breaking'', Section 6.1; independently formulated.)}
    {ch28-s14-the-gauge-mediation-problem}
Suppose the same one-loop messenger interaction generates both
  \(\mu\) and \(B\mu\).  Estimate their sizes in powers of
  \(F/S\) and a loop factor, then compare with the Higgs minimization
  conditions to exhibit the parametric tension.

  \SupplementaryExercise{15}{Radiative Breaking and Stable Cascade Endpoints}
    {(Inspired by Stephen P. Martin, \emph{A Supersymmetry Primer}, version 7, Sections 8.1 and 11.1; independently formulated.)}
    {ch28-s15-radiative-breaking-and-stable-cascade-endpoints}
  \begin{enumerate}
    \item \emph{Radiative Electroweak Breaking from the Top Yukawa.}
    Use the one-loop RGE for the up-type Higgs soft squared mass to explain
      how a positive high-scale value can run negative.  Track the competing
      top-Yukawa and gauge contributions with the correct direction of RG
      evolution.
    \item \emph{LSP Stability and Cascade Endpoints.}
    Prove from multiplicative \(R\) parity that sparticles must be produced
      in pairs, every heavier sparticle cascade ends in an odd particle, and
      the LSP is stable.  State which conclusions fail if a single
      lepton-number violating coupling is restored.
  \end{enumerate}

\end{enumerate}
